\documentclass{minmal}
\begin{document}
\textbf{Oppsummering til prøve om separering og blanding:}\nl[2]

\noindent I dette dokumentet angir jeg hva som er viktig for dere å lære til prøven.  Dette er snakk om stikkord som kan hjelp dere til å legge vekt på det jeg mener er det viktigste av stoffet.
\textbf{Blanding:}
Blandinger deler vi inn i  to typer: \textbf{homogene} blandinger og \textbf{heterogene} blandinger.   Dere må kunne definere disse (forklare de er) og gi et eksempel på hver av de.
\textbf{Suspensjon}: En suspensjon er en faststoff-væskeblanding.  Veldig ofte ønsker vi å kunne skille disse partiklene fra væsken.  Her vil det være slik at de partiklene som har mindre massetetthet en væsken (ofte vann) vil flyte opp (\textbf{flotasjon}). De som har større massetetthet enn vann  vil synke (\textbf{sedimentasjon}).  Noen partikler har massetetthet i nærheten av væsken.  Slike partikler kan vanskelig fjernes med sedimentasjon (tyngdekraft) og flotasjon (oppdriftskraft).  Da må vi bruke andre alternativ som f.eks.  sykloner (sentrifugalkraft) sentrifuger (sentrifugalkraft eller filtrering (trykkforskjell, vakuum).  I disse tilfellene innfører vi en ny kraft for å få skilt ut faststoffet.

\textbf{Emulsjoner}
Dere må¨kunne forklare hva som menes med en olje-i-vann emulsjon og en vann-i-olje emulsjon.  Dere må også kunne gi et eksempel på hver av dem. 

\textbf{Emulgator:} Dere må kort kunne si noe om hva en emulgator er (hva er oppgaven til den)?\\
\textbf{Saltløsninger:}
Er homogene løsninger.  Et godt eksempel er $NaCl$.  Når et salt løses i vann vil vannmolekylene "dra" ionen far hverandre, det er det som vi si at et salt er oppløst:
\begin{figure}[h]
    \centering
    \includegraphics[width=0.5\linewidth]{opplost_salt.png}
    \caption{oppløsning av et salt}
    \label{fig:placeholder}
\end{figure}\\
Vi ser at det som skjer er at vannmolekylene legger seg med sin positive ende rundt klor ionene, og med sin negative ende rundt natrium ionene, og drar ionene fra hverandre .  For noen salt er det lett for vann å dra ionene fra hverandre (lettløselige salt, $NaCl$) .  For noen salt er det vanskelig for vann å dra ionene fra hverandre (tungtløselige salt, $AgCl$).
Dere må kunne forklare umettet, mettet og overmettet løsning.  For $NaCl$ så får man løst opp ca. 360g per liter ved $25^0C$. Har vi mindre enn dette sier vi løsningen er umettet. har vi "akkurat" 360g/L sier vi at løsningen er mettet.  Og har vi mer enn 360g/L sier vi løsningen er  overmettet.  
\textbf{Gassblandinger}
Gassblandinger er homogene.  Der bør vite at luft inneholder $78\%\ N_2$ og ca. $21\%\ O_2$ 
\textbf{blandemetoder og utstyr:}
Vi konsentrerer oss om væskeblanding og mekanisk røring.  Her må dere lære dere alle begrepene til en standard rørebeholder:
\begin{figure}[H]
    \centering
    \includegraphics[width=0.5\linewidth]{rorebeholder.png}
    \caption{standard rørebeholder}
    \label{fig:placeholder}
\end{figure}

Spør om noen av disse er uklare.  
\textbf{Viskositet:} Viskositet kan enkelt forklares som hvor seig en væske er. Sirup har høy viskositet, vann lav viskositet.\newpage
For høy-viskøse væsker brukes saktegående rører (Wendel rører)
For lav-viskøse væsker brukes hurtiggående rører (propell rører)
Eksempel på rørere:
\begin{figure}[H]
    \centering
    \includegraphics[width=0.75\linewidth]{rørere.png}
    \caption{Noen rørertyper}
    \label{fig:placeholder}
\end{figure}
\noindent \textbf{Mekanisk separasjon:}\\\\
\textbf{faststoff faststoff separasjon:}
Her er det viktig å skille mellom \textbf{sortering} og \textbf{klassering}\\

\textbf{Sortering} handler om å skille to (eller flere) stoffer i en blanding etter type stoff (eks. skille epler og pærer).  
\noindent
\textbf{Klassering} handler om å separere stoffer etter størrelse. Størrelsen som er av betydning er diameteren på stoffene.  

NB! I dagligtalen kan vi ofte snakke om sortering når  vi deler inn etter størrelse. Men det gjør vi \textbf{ikke} i kjemisk teknologi.
\textbf{Utstyr til mekanisk separasjon av faste stoffer (sortering):}\\
\textbf{Sedimentasjonstank:\\}
Denne  ser slik ut:
\begin{figure}[H]
    \centering
    \includegraphics[width=0.5\linewidth]{sedimentasjonstank.png}
    \caption{sedimentasjonstank}
    \label{fig:placeholder}
\end{figure}
Til denne kommer det inn en suspensjon (hva var det?) som har faststoffpartikler med tetthet større enn vann og faststoffpartikler med tetthet mindre enn vann. De med mindre tetthet enn vann vil flottere (flotasjon) og de med tetthet større en vann vil synke (sedimentere).  Væsken blir satt i bevegelse av en \textbf{pulsator} eller et stempel som beveger seg opp og ned.
\textbf{Hydrosyklon:}\\
En hydrosyklon betyr at det er vann tilstede.  Man tilfører et vandig slam med partikler som har ulik tetthet:
\begin{figure}[H]
    \centering
    \includegraphics[width=0.5\linewidth]{hydrosyklon.png}
    \caption{hydrosyklon}
    \label{fig:placeholder}
\end{figure}
Her vil det være sånn at sentrifugalkraften er størst for partiklene med størst tetthet.  Disse  vil treffe syklonveggen og falle ned.  Partiklene med minst tetthet vil aldri treffe veggen. De vil gå i en indre spiral og ut på toppen sammen med vannet, se figuren over.\\
\textbf{Magnetseparator:}\\
Prinsippet for en magnetseparator er vist i figuren under:\\
\begin{figure}[H]
    \centering
    \includegraphics[width=0.4\linewidth]{magnetseparator.png}
    \caption{mangetseparator}
    \label{fig:placeholder}
\end{figure}
Her har vi en blanding av stoff som er magnetiske og ikke magnetiske.  P.g.a. den faststående magneten klarer vi å skille stoffene som vist i figuren. Spør om noe er uklart.\\\\

\textbf{Utstyr til mekanisk separasjon av faste stoffer (klassering):}\\
Når vi klasserer faststoff etter størrelse så er det partiklenes diameter vi går etter.  Dette gjøres på Saugbrugs for å klassere flis etter ulik størrelse- Hver størrelsesgruppe kalles for en \textbf{fraksjon}:
\begin{figure}[H]
    \centering
    \includegraphics[width=0.5\linewidth]{fraksjoner.png}
    \caption{ved sortering deler vi opp i ulike farksjoner}
    \label{fig:placeholder}
\end{figure}
\textbf{Utstyr til klassering:}\\
Som eksempel kan vi nevne trommelsikt:
\begin{figure}[H]
    \centering
    \includegraphics[width=0.5\linewidth]{trommelsikt.png}
    \caption{Enter Caption}
    \label{fig:placeholder}
\end{figure}
Denne burde vel være selvforklarende? Hvis ikke spør.\\

\textbf{Mekanisk separasjon av faststoff – væskeblandinger}

Vi tar for oss de vanligste metodene og utstyr som egner seg. Det fins helt sikkert mange flere både metoder og varianter av utstyr. Dette gir likevel en innføring i prinsipper for utstyr som er mye brukt.\newpage

Vi deler metodene inn i fire grupper av prosesser:

\begin{itemize}
    \item Sedimentering;
    \item Filtrering;
    \item Sentrifugering;
\end{itemize}
\textbf{Sedimentering}

Dette kalles også for sedimentasjon eller avsetning. Det er tyngdekraften som blir utnyttet. Tyngdekraften til faststoff partiklene i væskeblandingen må være større enn oppdriftskraften væsken har:
\begin{figure}[H]
    \centering
    \includegraphics[width=0.5\linewidth]{tyngdekraft-oppdriftskraft.png}
    \caption{Enter Caption}
    \label{fig:placeholder}
\end{figure}
 For partikler i en væske er det slik at:\\\\
 
 Når partikkelens massetetthet er mindre enn massetettheten til vann så vil partikkelen ha en oppdritskraft som er større en tyngdekraften.  Partikkelen vil dermed flyte opp (flottere)\\\\
 
 Når partikkelens massetetthet er større enn massetettheten til vann så vil partikkelen ha en tyngdekraft som er større en oppdriftskraften.  Partikkelen vil dermed synke til bunns (sedimentere)\newpage

 \textbf{Utstyr til sedimentering og klaring}

Siden fortykkere eller klaringsbasseng er veldig vanlig er det naturlig å ta med litt ekstra om denne. Her tar vi med en type, selv om det vil være annet utstyr som bruker de samme prinsippene.

\textit{Bildet viser prinsipp for en rund fortykker (klaringstank):}\\
\begin{figure}[H]
    \centering
    \includegraphics[width=0.75\linewidth]{rundt_klaringsbasseng.png}
    \caption{Et rundt klaringsbasseng}
    \label{fig:placeholder}
\end{figure}
Runde klaringsbasseng eller fortykkere er laget i ulike dimensjoner ut fra de mengdene væske / suspensjon de skal håndtere.  Her vil faststoffet synke til bunns mens de rensede vannet vil forlate tanken i overløp.  I bunnen er det en veldig saktegående skrape som vi skrape faststoffet inn mot midten og tas ut.\newpage

\textbf{Flotasjonsfortykker}

Prinsippet er at noe av faststoffet har lavere tetthet enn et annet fast stoff. Ved flotasjonsfortykker får vi skilt faststoff fra hverandre og fra væsken (vannet) som faststoffet befinner seg i.

\textit{Bildet viser separasjon av faststoff og væske i flotasjonsfortykker.}\\\\

\begin{figure}[H]
    \centering
    \includegraphics[width=0.75\linewidth]{flotasjonsfortykker.png}
    \caption{flotasjonsfortykker}
    \label{fig:placeholder}
\end{figure}
Flotasjon handler om å få noe til å flyte opp. Fortykning handler om å gjøre noe tjukkere og med større tetthet og størrelse. Da vil vi raskere få sedimentert ut faststoff med størst tetthet.

I tanken blir suspensjon (væske med faststoff) ført inn i toppen av tanken. Suspensjonen blir ledet inn i et «skjørt», slik at væske får retning nedover. Det gir raskere sedimentering av stoffer med høyest tetthet. Med litt tid vil tungt faststoff bli tatt ut som slam i bunnen av tanken. De letteste faststoff partiklene vil flyte med væsken opp legger seg som et lag over væsken. Dette blir tatt ut i et overløp eller skrapet av.

Væske som er fri fra både tunge og lette faststoff partikler blir tatt ut i et utløp mellom topp og bunn av tanken. I noen anlegg vil det som flyter øverst bli skrapt av med en rakearm. Rakearmen leder faststoffet mot et uttak. Utstyret eller metoden er mye brukt i avløpsrenseanlegg der for eksempel fettstoffer vil flyte opp, mens tyngre stoffer ikke vil gjøre det.

 \textbf{Filtrering}

Filtrering er en mekanisk separasjons metode til å skille faststoff fra en suspensjon. Vi får produktene filtrat (væsken) og fuktig faststoff. I noen tilfeller er det væsken vi er på jakt etter, – i andre tilfeller er det det faste stoffet. Separasjonen er mulig ved å bruke ulike typer filterelement.

Filterelementer finnes i naturen der vann renner gjennom stein og grus. I prosess sammenheng legger vi til rette for at filtreringen skal gi best mulig resultat. Filtreringen skjer ved at faststoff blir holdt tilbake av «et hinder», som for eksempel et filter, mens væske vil passere. Dermed kan vi skille faststoff som er i en væske. Prinsippet handler også om størrelse. Her trenger vi å vite diameteren på faststoffpartiklene, når vi skal bestemme diameter i åpningene i filterelementet.  Vi skal se nærmere på kakefiltrering.\\
\textbf{Kakefiltrering}

I prosessindustrien er det denne typen filtrering som er mest vanlig. Derfor skal vi bruke litt mer tid på denne typen filtrering.\\

\textit{Bildet viser prinsipp for kakefiltrering:}
\begin{figure}[H]
    \centering
    \includegraphics[width=0.5\linewidth]{kakefiltrering.png}
    \caption{Prinsippet for kakefiltrering}
    \label{fig:placeholder}
\end{figure}

Ved kakefiltrering blir det som oftest brukt en filterduk som kan være trukket over en trommel, et vandrebelte eller i rammer. Dette gir ulike typer filtreringsapparater som vi skal se nærmere på.

Kakefiltreringen går ut på at maskeåpningene i filterduken er mindre enn diameteren til faststoffpartiklene vi ønsker å skille ut. I starten av prosessen vil de minste partiklene gå igjennom, men etter hvert som faststoff samler seg på filteret, fungerer også laget med faststoff som et filter. Vi kaller faststoffet som ligger utpå filteret for filterkake.

 
\textbf{Vakuum trommelfilter}

Utstyret består av en stor trommel med huller i. Den kan minne om trommelen i en vaskemaskin. Trommelen blir trekt med filterduk. Inni trommelen er det ulike soner, og under trommelen finner vi et trau. Trauet er en beholder for suspensjonen som skal filtreres. I noen trau er det omrører, som for eksempel en propell. Andre trau skaper bevegelse ved at den bli vugget. Disse bevegelsene – eller omrøringen er viktig for at faststoffet i suspensjonen ikke skal sedimentere til bunn (setle seg).\newpage

\textit{Bildet viser prinsipp for vakuum trommelfilteret :}\\\\
\begin{figure}[H]
    \centering
    \includegraphics[width=0.75\linewidth]{trommelfilter.png}
    \caption{trommelfilter}
    \label{fig:placeholder}
\end{figure}

Suspensjonen blir altså tilført i trauet. En del av trommelen er nedsenket i suspensjonen og roterer sakte rundt.

Inni trommelen er det undertrykk (vakuum) som blir produsert av en sentrifugalpumpe. Suget gjør at suspensjon trekkes opp mot filteret på den delen som er nedi suspensjonen. Siden faststoffet har større diameter enn maskeåpningene i filterduken, blir faststoff partiklene hengenes igjen på utsiden av duken. Væsken (filtratet) går igjennom filteret og sendt i retur etter rensing. Renset filtrat bidrar da med vask av filterkaken etter hvert som trommelen roterer.

Når trommelen kommer rundt blir filterkaken fjernet. Dette kan skje ved at kniver kutter av kaken, slik at filterkake detter av. Alternativt kan siste sektor av trommelen kobles til trykkluft i stedet for vakuum, slik at kaken blir blåst av. Filterkaken transporteres videre enten med et belte eller på annet vis. I denne prosessen er det filterkaken som er det viktige produktet.

For at prosessen skal gå kontinuerlig er det viktig å holde filteret ved like, eller bytte det ut. Dersom det blir målt mye fuktighet i filteret kan det skyldes at filteret har huller. Filterkaken skal være passe tykk og passe tett. Blir den for tett kan ikke vannet trenge forbi. Blir den for tynn kan det skyldes at pumpa ikke skaper nok undertrykk til at filteret kan trekke til seg mer suspensjon.\newpage

 \textbf{Faststoff-gass separasjon:}\\\\
 Mye av utstyret og metodene som hører til separasjon av faststoff-gass vil du finne igjen i gassrenseanlegg. I forhold til tema miljø er det et viktig ledd å fjerne støv fra en gass før den slippes ut til luften. Å skille faststoff fra gass kan likevel ha andre positive sider. For bedriften sin del handler det om økonomi. Faststoffet kan inneholde råstoffer til prosessen. Ved å hente dette ut fra avgasser kan bedriften utnytte mer av råstoffet eller gjenvinne stoff som er på vei til å «gå tapt». I tillegg vil bedriften få økonomiske konsekvenser dersom utslipp av støv til luft overstiger konsesjoner de har fått fra myndighetene.

I denne delen ser vi på mekaniske metoder til å separere faststoff i en gass, i tillegg til elektrostatiske metoder og noen pneumatiske metoder. Vi kan dele mekanisk støvfjerning i tørr støvfjerning, våt støvfjerning og støvfjerning ved bruk tørrfilter eller elektrofilter.

 \textbf{S}\textbf{Støvkammer}

Denne virker etter prinsipp om \textit{tyngdekraft}. Støvkammeret er en stor beholder med et smalt inntak for gass som skal renses, en innsnevring i bunn med uttak for fast stoff som skal tas ut og et smalt uttak for renset gass.
\textit{Bildet viser et støvkammer med avbøyningsplater .}
\begin{figure}[H]
    \centering
    \includegraphics[width=0.75\linewidth]{stovkammer.png}
    \caption{Støvkammer}
    \label{fig:placeholder}
\end{figure}

Rågassen, gassen som skal renses, inneholder faststoff partikler. Rågassen blir tilført kammeret gjennom den trange passasjen. Vi er kjent med at når vi har liten diameter (tverrsnitt) på en åpning vil hastigheten være høy og trykket lavt, mens når diameteren øker vil hastigheten falle og trykket øke. Dette skjer også for gassen som kommer inn i det store kammeret. Med lavere hastighet får tyngdekraften \textit{tid }til å virke på faststoff partiklene slik at disse vil falle mot bunnen av kammeret.

 
\textbf{Tørr syklon}

Sykloner er veldig vanlige utstyr i prosesser og anlegg. Tørr Syklon behandler tørr gass, med mål om å skille ut faststoff partikler som følger med rågassen. Prinsippet i syklon er \textit{sentrifugalkraft }og \textit{tyngdekraft}. Fast stoff har større masse og tetthet enn gass. Det er dette som gjør det mulig å skille ut faststoff partiklene.:\\\\
\begin{figure}[H]
    \centering
    \includegraphics[width=0.5\linewidth]{torrsyklon.png}
    \caption{tørrsyklon}
    \label{fig:placeholder}
\end{figure}
 Faststoffpartiklene vil treffe veggen og falle nedover langs veggen. Da tas de ut i bunnen.  Men gassen vil aldri treffe veggen. Den vil gå i en indre spiral og ut på toppen.\newpage
 \textbf{Posefilter}

Posefilter har prinsipp som støvsuger. I støvsugeren har vi poser til oppsamling av støv. Når vi starter støvsugeren vil motoren bidra til at det skapes sug som trekker luft og støv inn i posen. Støvet blir værende igjen i posen, mens luft går ut via et filter i toppen. Når posen er full suger støvsugeren dårlig. Da må vi skifte filterposen.  \textit{Bildet viser forenklet hvordan posefilter er bygget opp og fungerer}. \\\\
\begin{figure}[H]
    \centering
    \includegraphics[width=0.75\linewidth]{posefilter.png}
    \caption{posefilter}
    \label{fig:placeholder}
\end{figure}\newpage
\textbf{Elektrofilter}

Elektrofilter kombinerer prinsipper om at støv/fast stoff kan få en elektrisk ladning når de utsettes for strøm, mens gasser ikke lar seg lade, – og at faststoff har en masse og tyngde som gjør at tyngdekreftene vil virke på faststoffet.  \textit{Bildet viser prinsipper som skjer ved elektrofiltrering.}:\\\\
\begin{figure}[H]
    \centering
    \includegraphics[width=0.5\linewidth]{forklaring_el_filter.png}
    \caption{Prinsippet for elektrofilter}
    \label{fig:placeholder}
\end{figure}

Rågass som skal renses for faststoff partikler, blir bombardert av elektroner fra en negativt ladet elektrode. Veggene i elektrofilteret er koblet mot positiv ladning. Støv som har blitt negativt ladd blir derfor med til den positivt ladde veggen. Her vil støvet bli utladet, siden pluss og minus blir nøytralt. Etter hvert samler det seg mer støv på veggen og tyngdekraften vil få støvet til å falle mot bunn.\newpage  
\textit{Bildet viser prinsipp for plateelektrofilter:} \\\\

\begin{figure}[H]
    \centering
    \includegraphics[width=0.5\linewidth]{plate-elektrofilter.png}
    \caption{olate-elektrofilter}
    \label{fig:placeholder}
\end{figure}
\end{document}